\section{논의}
\label{sec:discussion}

%======================================================================
\subsection{단계 위험을 좌우하는 요인}
\label{sec:discussion:drivers}

% P1: 학습 목적이 추정 대상을 정한다
\rev{학습 목적은 모델이 무엇을 추정할지를 정하며, 다음 단계 탈락 위험에서는 이 선택이 설계 간 차이의 대부분을 만들었다.
상대위험 목적은 단계마다 공통적인 위험 수준을 배우지 못해 예측이 평평해지지만(\ref{sec:method:rq2}절, 그림~\ref{fig:contrast}(a)), 목적을 바꾼 이득은 재보정의 영향을 받지 않는 단계 안·경기 안 비교에서도 남았다.}
이는 다음 단계라는 짧은 구간에 맞춘 목적이 남은 경기 전체의 탈락 순서를 맞추는 목적과 다른 정보에 무게를 두기 때문으로 해석된다.
\rev{반대로 KRAFTON 적용판과 Cox 목적은 원래 과제에서 앞섰으므로(\ref{sec:results:rq2}절), 설계의 비교는 성능 순위보다 추정 대상이 의사결정의 시간 단위와 맞는지를 기준으로 읽어야 한다.}

% P2: 입력은 구역과 주변 팀에 대한 위치
\rev{입력에서는 특징의 수보다 줄어드는 구역과 주변 팀에 대한 팀의 위치를 담았는지가 중요하였다.
그런 정보가 없는 활동 점수에 서식지 맥락 블록을 더하면 Sagong 설계와의 차이의 절반 이상이 메워졌고, 본 연구의 입력에서 이 블록을 빼면 성능이 떨어졌다.}
구역까지의 거리와 주변 팀 수를 이미 담은 KRAFTON 입력은 본 연구의 입력과 예측력이 같았지만 둘을 합치면 성능이 올랐는데, 무기별 피해와 장비 수준은 KRAFTON 입력에만, 이동 부담과 관계 요약은 본 연구의 입력에만 있다는 점에서 두 입력은 서로 다른 측면을 담는 것으로 보인다.

% P3: 구조는 경기 간 맥락에 작게 기여
\rev{모델 구조는 같은 경기 안의 순위를 바꾸지 않았고, 같은 단계에서 경기마다 다른 위험 수준을 담는 데 작게 기여하였다.
팀 사이의 관계는 주변 경쟁 밀도와 관계 요약 특징으로 이미 입력에 요약되어 있어 같은 경기 안에서 attention이 새로 찾을 정보는 적었던 반면, 팀 축은 함께 살아 있는 팀 전체를 보므로 경기 단위의 맥락을 더할 수 있었던 것으로 해석된다.
이 차이는 통합 AUC만으로는 보이지 않는다.}

% P4: 생태학적 해석
서식지 관점은 \rev{요인의 연관이} 단계에 따라 달라지는 양상을 해석하는 틀이 된다.
구역이 넓은 초반에는 주변 팀의 수가 위험을 크게 가르지만, 구역이 좁아진 후반에는 모든 팀이 가까이 모여 그 차이가 덜 중요해지고 남은 구역 안의 위치와 남은 팀원 수가 더 중요해지는 것으로 해석된다.
\rev{경계에 가까울수록 위험이 높고 그 차이가 구역이 작아질수록 커지는 양상은, 서식지 가장자리의 개체가 더 큰 위험에 놓이는 가장자리 효과와 비슷하다~\cite{Ries04}.
이는 관측된 연관에 대한 해석이지만, 네 변수가 같은 단계 안에서도 판별력의 약 7할을 재현했다는 점에서 이 틀은 위험을 요약하는 데 쓸모가 있다.}

%======================================================================
\subsection{활용과 시사점}
\label{sec:discussion:implications}

% P5: 경기 분석, 중계, 구역 설계
\rev{단계별 평균 위험 수준이 맞는 예측은 단계 사이의 위험을 같은 기준점에서 비교하게 하므로, 경기 후 복기, 중계 분석, 구역 규칙 변경 전후의 비교에 쓰일 수 있다.
다만 실시간 중계에 쓰려면 실시간 데이터 공급과 위험 구간별 보정의 확인이 더 필요하다.}

% P6: 대규모 시공간 위험 예측에 대한 시사점
이 결과는 이동하는 다수의 에이전트를 다루는 대규모 시공간 위험 예측에도 세 가지 시사점을 준다.
첫째, 입력뿐 아니라 생존 판정과 에이전트 간 관계에도 같은 절단 시각을 적용하고, 절단 이후 기록을 지운 재구성으로 누수가 없음을 확인해야 한다.
둘째, 통합 AUC는 시점별 위험 수준을 맞히는 효과와 같은 상황의 개체를 가려내는 효과를 섞으므로, \rev{설계 요인의 효과도 단계 안·경기 안 지표와 단계별 보정으로 나누어 보고해야 한다.}
셋째, 설계를 비교할 때는 구조를 바꾸기 전에 추정 대상과 입력을 통제한 분해로 차이의 출처부터 확인해야 한다.

%======================================================================
\subsection{한계}
\label{sec:discussion:limitations}

% P7: 한계
본 연구의 결론은 다음 범위 안에서 읽어야 한다.
첫째, 입력은 모든 팀의 위치를 아는 관전자 관점에서 만들었고 \rev{예측 시점에 공개되는 다음 구역은 쓰지 않았으므로}, 각 팀이 당시 알던 정보\rev{로 하는} 의사결정 지원과는 구분된다.
둘째, 자료는 한 구역 규칙 아래 약 10주간의 스팀 서버 경기를 \rev{경기 단위로 무작위 분할한 것이고, 솔로 코호트는 하나뿐이며 일부 코호트에는 AI 계정이 많다. 다른 규칙·플랫폼·시기와 솔로 경기에 일반화하려면} 추가 확인이 필요하다.
셋째, 두 기존 설계는 원래 자료와 조건을 재현한 것이 아니라 같은 자료에 맞춘 적용판이며, \rev{특히 KRAFTON 설계는 16팀 이하의 대회 경기를 대상으로 만들어졌고 일부 입력은 근사하였으므로} 원 논문의 수치와 직접 비교할 수 없다.
넷째, RQ1의 계수는 관측 자료의 조건부 연관이며, \rev{팀이 위치를 스스로 고르고 탈락 원인(구역 피해와 교전)을 구분하지 않았으므로 인과 효과로 읽을 수 없다.}
다섯째, \rev{구조에 관한 결론은 현재의 데이터 규모와 요약 특징, 코호트마다 하나의 초기화로 학습한 조건의 결과이며}, 더 많은 경기나 플레이어 수준의 궤적을 쓰면 달라질 수 있다.

%======================================================================
\section{결론}
\label{sec:conclusion}

% P8: 결론과 향후 과제
본 연구는 배틀로얄 경기의 대규모 텔레메트리에서 다음 단계 탈락 위험을 서식지 축소 관점의 단계 조건부 이산시간 hazard로 정식화하고, 누수 없는 팀--단계 표본과 통제된 분해로 이 위험의 예측을 무엇이 좌우하는지 분석하였다.
경쟁 밀도, 이동 부담, 남은 팀원 수의 \rev{연관은} 코호트에 걸쳐 일관되었고 경쟁 밀도와 구역 안쪽 깊이의 \rev{연관은} 단계에 따라 달라졌다.
\rev{기존 설계와의 차이는 주로 정식화와 서식지 맥락에서 왔고, 모델 구조는 같은 경기 안의 순위가 아니라 경기 간 위험 차이에 작게 기여하였다.}
향후에는 각 팀이 관측할 수 있는 정보만 쓰는 의사결정 관점으로 모델을 확장하고, 수집 이후 적용된 구역 규칙 개편~\cite{PUBGBlueZone2026}을 자연 실험으로 삼아 위험 요인이 규칙이 바뀌어도 유지되는지 검증할 계획이다.

\section*{사사}
\rev{Anthropic의 생성형 AI인 Claude(Opus 5.5)를 분석·그림 코드의 작성과 초록 및 I--VI절의 초안 작성·교정에 사용하였다.
저자는 AI가 만든 모든 코드, 수치, 문장을 검토하고 수정하였으며 그 내용에 책임을 진다.}
